\chapter{Reuseable writing}

\paragraph{Please write a Statement of Interest/Personal Statement explaining your goals, experiences, and skills as it relates to this program.}

My primary motivation is the intellectual challenge of pushing beyond the boundaries of current knowledge. There is always more to learn, more fundamental principles to understand, and more open problems to solve. What draws me to research is that it offers a continuous challenge: when one problem is solved, new questions emerge, and when original research is not possible, there is still the opportunity to expand one's understanding through entirely different frameworks and formalism's rather than just simply adding knowledge to an already constructed base. I find particular value in learning abstract concepts from first principles, building an understanding from fundamental axioms.

Physics is especially appealing to me because of its deep connection with mathematics. The same physical phenomenon can often be understood through multiple mathematical formalisms, each revealing different structures, assumptions, and insights. Exploring these perspectives provides an almost limitless opportunity to learn, question established ideas, and develop new ways of approaching difficult problems. For me, the appeal of physics lies not only in understanding how the world works, but also in understanding why our mathematical descriptions work and how they can be extended.

I hope to make this intellectual pursuit a career through research in plasma physics at a national laboratory, particuraly through the computational and theoretical plasma physics programme. I want to work on challenging, consequential problems that require both mathematical depth and computational ingenuity, drawing on different theoretical frameworks and formalisms to develop new approaches.

This program will let me engage in these problems, through fusion and energy research, developing the skills I need to push further than I already have. Oak Ridge being at the forfront of this research area, the researchers there would be able to teach me what is important at speeds exponentially faster than my self-education stumbling through individual papers.

As for my ability to help, doing research with Dr. Mezzacappa has taught me a substantial amount about signal theory, magnetohydrodynamics, kinetic theory, numerical methods, and scientific computing. While we used signal theory for linear match filtering of gravitational wave data, it is also extremely useful in diagnosing the extremely noising enviorment due to the extremely instabilities caused by plasma interactions. While our equation is used for the interactions inside of core collapse supernova of plasma in fluid flow as well as the Boltzmann equation for neutrino interactions, fusion physics is almost entirely built upon MHD as well as extending the Boltzmann equation into Vlasov and gyrokinetic equations. Finally, the usage of FORTRAN, MPI, supercomputers, and data analysis using python lays the foundation of fusion as the plasma interactions are too complex to be solved analytically.

Futher, I have done substantial independent projects that push my understanding further. I have a YouTube channel where I have built numerical simulations using discontinuous Glarklin, derived Boltzann equation from N-body Hamiltonian and Navier-stokes from Boltzmann, PPO AI playing an asteriod game, integrate Resistive Magnetohydrodyanmics equation using variational integrator, General Relativistic BSSN numerical code, and examining Metriplectic dynamics. Most of these are examining numerical methods useful in plasma dynamics. Other develop the mathematical formulisms very useful to developing plasma physics.

\paragraph{What about SMART interest you the most}


\paragraph{Interest in Undergraduate research}


My primary motivation is the intellectual challenge of pushing beyond the boundaries of current knowledge. There is always more to learn, more fundamental principles to understand, and more open problems to solve. What draws me to research is that it offers a continuous challenge: when one problem is solved, new questions emerge, and when original research is not possible, there is still the opportunity to expand one's understanding. I find particular value in learning abstract concepts from first principles, building an understanding from fundamental axioms rather than simply accepting established results.

Physics is especially appealing to me because of its deep connection with mathematics. The same physical phenomenon can often be understood through multiple mathematical formalisms, each revealing different structures, assumptions, and insights. Exploring these perspectives provides an almost limitless opportunity to learn, question established ideas, and develop new ways of approaching difficult problems. For me, the appeal of physics lies not only in understanding how the world works, but also in understanding why our mathematical descriptions work and how they can be extended.

I hope to make this intellectual pursuit a career through research in plasma physics at a national laboratory. I want to work on challenging, consequential problems that require both mathematical depth and computational ingenuity, drawing on different theoretical frameworks and formalisms to develop new approaches. Ultimately, I want a career in which learning and discovery are inseparable from one another: one that continually challenges me to deepen my understanding while contributing original work toward solving difficult scientific problems.
\paragraph{Why this project and this mentor}

Dr. Mezzacappa is working on Core-Collapse supernova, specifically we are currently working on developing ways to detect
gravitational wave memory with next generation detectors.

This is extremely exciting work as it lays to groundwork to eventually verify another validation of General Relativity or find an
interesting problem in Relativity that asks extremely challenging problems with it.

Beyond that, the underlying math of simplifying the extremely hyperbolic equations of GR and the highly non-linear equations
of CCSN, leads to relatively simple computations to calculate gravitational waves to surprisingly high accuracy.

Further, I have done work with Dr. Mezzacappa during high school, and I know him to be extremely kind, altruistic, and he actively looks out for my well being and my future which I am incredibly thankful for.

\paragraph{What I want to learn}

Most directly I am already learning a lot about signal theory and data analysis, extremely important tools in all of science and
pretty far beyond.

I am also observing how science happens; how they talk about these very abstract topics in concrete ways to intuitively figure
out what directions to go, what is feasible, and what is useful in a topic that is outside the scope of modern human knowledge.
Also, before and after discussions I get to ask questions on tangentially related topics to learn much more about the vast
physics surrounding stars and their respective supernova.
